\documentclass[twoside]{article}
\hfuzz=\maxdimen
\usepackage[
    a5paper,
    margin=0.5in,
    includehead,
    headsep=4pt,
    headheight=13pt
]{geometry}
\usepackage{gregosheet}
\usepackage{gregosheet-psalm}
\input{styles}
\input{ordo-cantus-officii/1-proprium-de-tempore.tex}
\input{ordo-cantus-officii/2-psalterium-per-quattuor-hebdomadas-distributum.tex}

% ============================================================================
% IULIUS
% ============================================================================

% ············································································
% Die 11: S. Benedicti
% ············································································

% --- Officium Lectionis ---

% Hymnus
\defsheet{InterAeternas}
  {<X-5-4--5--4----1í--:-4--4---3---4--5x--5x-,-5---7--8----7---5x-:-5--4-3---4--5x-5x-,-5---4--5-------4---1í-:-4--4--3---4--tT4-3y-,-3--3--2--1í-1í-?,-1wW1-0\\q\\-.}
  {   É-gi fé-nyek közt, ki-ket áld-va hir-det  har-ci győz-tes-ként di-a-dal-mi é--nek, ott ra-gyogsz, bol-dog  Be-ne-dek su-gár-zó   ér-de-me-id-del.  Á----men.}

% --- Laudes ---

% Hymnus
\defsheet{LegiferPrudens}
  {<X-5-----4--5--4-----1í--:-4--4--3----4--5x-5x-,--5-----7--8----7--5x-:-5-4--3---4--5x-----5x--,-5---4--5--4---1í-:-4-4--3----4--tT4--3y-,-3--------3--2---1í--1í-?,-1wW1-0\\q\\-.}
  {   Bölcs ta-ní-tónk, Szent Be-ne-dek, vi-lá-gíts, Krisz-tu-sunk fé-nyét a vi-lág-ba hintsd szét, hí-ven is-mer-nünk e hi-tet, se-gíts meg, s_töltsd be szí-vün-ket!  Á----men.}

% Benedictus
\defsheet[tone={8. dó tónus}]{FuitVir}
  {<-7-5u-7--¨-7iI7-uU6u-uU6-5--6u--5--4---,-rR3-4--3--5--uU67iI7-:-5u-6---5---4---3r-4-?,-¦-5-7-8-7-.}
  {  É-le-te * tisz-te---let-re-mél-tó volt, ne--ve Be-ne-dek:      ke-gye-lem-mel ál-dott.}

% ============================================================================
% AUGUSTUS
% ============================================================================

% ············································································
% Die 8: S. Dominici
% ············································································

% --- Laudes ---

% Hymnus
\defsheet{NovusAthletaDomini}
  {<ô-2--2---2--1---2---4---3--2í-:-2--3---4----5--4---5---5--6x-,-6--6--5---6---4--5--rR3-2í-:-2--1----2-4--5---2---1--2í-?,-1e4R3E2W1-2í-.}
  {   Di-cső ne-vet mél-tán vi-selt Do-mon-kos, Is-ten baj-no-ka,  az Úr-hoz fűz-te őt ne--ve,  és lett a jó hír har-co-sa.   Á---------men.}

% Benedictus
\defsheet[tone={1. ré tónus}]{Beatus}
  {<-0q--1tz-5--5z-4tT4---:-4-5--6--5--rR3-3rR3-1---:-4--5--4--3---eE2--1---2---0-,-0--1--3-:-3-----3--4t--wW1-3--2---1--1-?,-[-4-3-4t-4-.}
  {  Bol-dog az a  szent, * a-ki az Úr-ban-bí---zott, az Úr pa-ran-csát hir-det-te, és az ő   szent he-gyé-re  ál-lít-ta-tott.}


\input{ordo-cantus-officii/4-communia.tex}
\input{ordo-cantus-officii/6-appendix.tex}

\title{Tábori lectionarium}
\subtitle{902. Kucsera Ferenc cserkészcsapat - 2026. nyári tábor}

\begin{document}

\maketitle

\section*{Évközi 14. vasárnap - Előtábor 2. nap}

\parttitle[12, 1-25]{Első olvasmány}
Sámuel második könyvéből

\begin{lectio}{Dávid bűnbánata}
\noindent Azokban a napokban az Úr elküldte Nátán prófétát Dávidhoz. El is ment hozzá, és így
beszélt: „Egy városban élt két ember. Az egyik gazdag volt, a másik szegény. A gazdagnak rengeteg
juha, kecskéje és barma volt. A szegénynek azonban nem volt egyebe, csak egy kisbáránya, vette
magának. Etette, úgyhogy megnőtt, vele magával és gyermekeivel. Evett a kenyeréből, és ivott a
poharából. Az ölében aludt, s olyan volt, mintha a lánya lett volna. Egyszer vendég érkezett a
gazdag emberhez. De nem tudta rászánni magát, hogy a juhai vagy barmai közül egyet is elvegyen, és a
vendégnek, aki érkezett, elkészítse. Ezért elvette a szegény ember báránykáját, és azt készítette el
vendégének.” Dávid nagyon megharagudott az emberre, és azt mondta Nátánnak: „Amint igaz, hogy az Úr
él: az az ember, aki ilyet tesz, az a halál fia. A bárányt négyszeresen vissza kell fizetnie, amiért
így tett, és nem volt kímélettel.”

Erre Nátán így szólt Dávidhoz: „Magad vagy ez az ember! Ezt mondja hát az Úr, Izrael Istene:
Fölkentelek Izrael királyává, és kiszabadítottalak Saul kezéből. Neked adtam urad házát és kebledre
urad asszonyait, rád bíztam Izrael házát és Júda házát. S ha ez kevés lett volna, még megtetéztem
volna ezzel meg azzal. Miért vetetted hát meg az Urat, s tettél olyat, ami gonosznak számít az Úr
szemében? Kard által elveszítetted a hetita Úriját, hogy feleségül vedd a feleségét. Igen,
elveszítetted az ammoniták kardja által! Nos, a kard sose fordul el házadtól, amiért megvetettél, és
elvetted a hetita Úrijától a feleségét, hogy a te asszonyod legyen. Ezt mondja az Úr: Nos, a saját
házadból hozok rád romlást. Elveszem a szemed láttára asszonyaidat, és odaadom őket más valakinek,
hogy napvilágnál asszonyaiddal háljon. Te titokban tetted, de én egész Izrael szeme láttára,
napvilágnál viszem ezt végbe!”

Erre Dávid így szólt Nátánhoz: „Vétkeztem az Úr ellen!” Nátán ezt válaszolta Dávidnak: „Így az Úr is
megbocsátja bűnödet, s nem halsz meg. De mert ezzel a tetteddel káromoltad az Urat, a fiú, aki
született neked, meghal.” Ezzel Nátán elment haza.

Az Úr megverte a gyermeket, akit Úrija felesége szült Dávidnak, úgyhogy súlyosan megbetegedett.
Dávid az Úrhoz fordult a fiú miatt; szigorú böjtöt tartott, s amikor hazament, éjszaka a puszta
földön aludt zsákkal takarózva. Udvarából a főemberek eléje járultak, és igyekeztek rábeszélni, hogy
keljen föl a földről, de nem akart, s nem is evett velük. A hetedik nap a gyermek meghalt. Dávid
udvari emberei nem merték neki jelenteni, hogy a gyerek meghalt. Azt mondták magukban: „Már akkor
sem hallgatott ránk, amikor még élt a gyerek, s a lelkére beszéltünk. Hát hogy mondhatnánk akkor meg
neki: a fiú halott? Még kárt tenne magában!” De Dávid észrevette, hogy udvari emberei suttogtak
egymás közt. Ebből megértette, hogy a gyerek meghalt. Ezért Dávid megkérdezte udvari embereitől:
„Meghalt a gyerek?” „Meg” – felelték.

Erre Dávid fölkelt a földről, megmosdott, bekente magát olajjal, és más ruhát öltött. Aztán elment
az Úr házába, és leborult. Ezután ételt tettek eléje, és evett. Erre udvari emberei így szóltak
hozzá: „Mit csinálsz? Amikor még élt a gyerek, böjtöltél, és siránkoztál. Most meg, hogy a gyerek
meghalt, fölkelsz, és ételt veszel magadhoz.” Ezt válaszolta: „Amíg a gyermek még élt, böjtöltem és
sírtam, mert azt gondoltam, ki tudja, hátha megkönyörül rajtam az Úr, és életben marad a gyerek. De
most, hogy meghalt, minek böjtöljek tovább? Hát visszahozhatom? Én utána megyek, de ő nem tér hozzám
vissza.”

Dávid Batsebát is vigasztalta. Elment hozzá, és együtt hált vele. (Batseba) fogant, és fiút szült. A
Salamon nevet adta neki. Az Úr kedvét találta benne, s ezt tudtul is adta Nátán próféta által, aki
az Úr szavára Jedidjának nevezte el.
\end{lectio}

\parttitle[Manassze imája 9. 10. 12; Zsolt 50, 5. 6]{Válaszos ének}

\versicle{%
    Megsokasodtak vétkeim, és bűneimnek száma, mint a teng\underline{e}r hom\underline{o}kja; †
    bűneim sokasága miatt nem vagyok méltó, hogy lássam a magas eget, mert haragra
    inger\underline{e}lt\underline{e}lek, * És ami színed előtt gonosz, \underline{o}lyat
    t\underline{e}ttem.
}{%
    Gonoszságomat bel\underline{á}tom, † bűnöm előttem lebeg szünt\underline{e}len, * egyedül csak
    ellened v\underline{é}tk\underline{e}ztem.
}

\parttitle{Második olvasmány}
Szent Ágoston beszédeiből

{\centering\redtext{(Serm. 19, 2-3: CCL 41, 252-254)}}

\begin{lectio}{Istennek tetsző áldozat a megtört szív}
\noindent \textit{Gonoszságomat belátom} (Zsolt 50, 5), mondja Dávid. Ha én belátom, akkor te ne
ródd föl azt nekem, Istenem. Még ha becsületesen éltünk is, ne gondoljuk, hogy nincs bűnünk. Életünk
annyira lesz dicséretes, amennyire Isten bocsánatát kérjük. Kétségbeejtő az olyan ember, aki nem
annyira saját bűneit nézi, de annál inkább vájkál a másokéban. Szimatolja, de nem azért, hogy
javítson rajta, hanem mert marni akarja. Mivel magát nem tudja tisztára mosni, ezért inkább másokat
igyekszik bemocskolni. Mi ne így tegyünk, amikor Dávid példájára imádkozunk, és elégtételt akarunk
nyújtani Istennek. Dávid azt mondta: \textit{Gonoszságomat belátom, bűnöm előttem lebeg szüntelen}
(Zsolt 50, 5). Nem figyelt ő más bűnére, hanem saját magát vádolta. Nemcsak a felszínen
tapogatódzott, hanem mélyen magába szállt. Saját magának nem kegyelmezett, hogy a maga számára ne
meggondolatlanul kérje a megkegyelmezést.

Ki akarsz engesztelődni Istennel? Gondolkozz: mit kell tenned, hogy Isten megbékéljen veled. Figyeld
meg, mit olvasunk ugyanebben a zsoltárban: \textit{Te a véres áldozatot nem kedveled,
égőáldozataimban nem leled kedved} (Zsolt 50, 18). Hát akkor ne is legyen már áldozatod? Ne ajánlj
fel neki semmit? Hát semmiféle áldozattal sem lehet Istent kiengesztelned? Hogy is mondtad a
zsoltárban? Te a véres áldozatot nem kedveled, égőáldozataimban nem leled kedved. Folytasd csak
tovább, és halld: \textit{Isten előtt kedves áldozat a megtört lélek, te nem veted meg a töredelmes
és alázatos szívet} (Zsolt 50, 19). Vesd el, amit hajdan áldoztál, hiszen megtaláltad, hogy ezután
mit áldozzál. Hajdan atyáid állatokat öltek le az oltáron, és azt tartották áldozatnak. Te a véres
áldozatot nem kedveled. Tehát azokat már ne keresd, hanem keresd a te áldozatodat.

Égőáldozataimban – mondja – nem leled kedved. Tehát ha égőáldozataimra rá sem nézel, vajon akkor
semmiféle áldozat sem jó neked? Dehogynem! Isten előtt kedves áldozat a megtört lélek, te nem veted
meg a töredelmes és alázatos szívet. Van tehát, mit áldozzál. Hát akkor ne nyájak körül nézelődj, de
hajóútra se készülj, hogy távoli országokból drága illatszereket hozzál. Saját szívedben keressed,
mi volna kedves Istennek. Szívedet kell megtörnöd. Mit félsz attól, hogy elvész, ha összetöröd? Hisz
ott áll a zsoltárban megírva: \textit{Tiszta szívet teremts bennem, Istenem!} (Zsolt 50, 12) Hogy
tiszta szívet lehessen benned teremteni, ahhoz a tisztátalant össze kell törnöd.

Ne legyen tetszésünkre az, ha vétkezünk, mert a bűn nem tetszik Istennek. Mivel nem vagyunk bűn
nélkül, legalább abban hasonlítsunk Istenhez, hogy nekünk se legyen a tetszésünkre az, ami neki nem
tetszik. Legalább annyira egyesülj Isten akaratával, hogy neked sem tetszik tenmagadban, amit utál
az, aki téged teremtett.
\end{lectio}

\parttitle[Zsolt 50, 12]{Válaszos ének}

\versicle{%
    Bűneim, Uram, mint a nyilak átszög\underline{e}ztek \underline{e}ngem, † de mielőtt sebeim
    elmérges\underline{e}d\-n\underline{é}nek, * Uram, Istenem gyógyíts meg engem a bűnbocsánat
    gy\underline{ó}gyszer\underline{é}vel.
}{%
    Tiszta szívet teremts bennem, Ist\underline{e}nem; * új és erős lelket \underline{ö}nts
    b\underline{e}lém.
}

\end{document}